% Einheit 2 von 4 – Parameter aus der Lagebedingung (bank/lagebeziehungen/e2.jsonl, zusammenbau v0.1)
%% TODO Zeitmarke Einheit 2: Marken-Zeile nicht lesbar
%% TODO Zweigzeile Teil 1: was hier gelernt wird (Satz in Schülersprache aus dem Einheitstitel) – Einheit 2
\einheitenkopf[e2]{Einheit 2 von 4 · Parameter aus der Lagebedingung}
\zweigzeile{Abitur LK}
\setcounter{aufgabe}{11}

%% TODO Titel: Ich-kann-Satz fehlt in der Bank (Platzhalter: Kettenname) – Nr. 12
%% TODO Anweisung: ein Satz über der ersten Teilaufgabe fehlt in der Bank – Nr. 12
\begin{aufgabe}{Parameter aus der Lagebedingung}
%% TODO \rechenplatz{n}: Schrittzahl des Grundfalls fehlt in der Bank (nur bei fester Schreibform, 2.3 b)
\begin{teile}
\teil Für welche Zahl $k$ liegt $P(1 | 2 | 0)$ in der Ebene $E_k: kx + 3y + z = 8$? Wer gegen wen? \\ \kreuz{Punkt gegen Ebene} \\ \kreuz{Gerade gegen Ebene} \\ \kreuz{Punkt gegen Gerade}
\teil $E_k: kx + 2y - z = 5$ enthält $P(2 | 1 | 3)$. Bestimme $k$. $k = $ \leerfeld
\teil Das Dreieck mit $A(3 | 0 | 0)$, $B(0 | 6 | 0)$ und $C(0 | 0 | 2)$ liegt in einer Ebene der Form $2x + y + kz = 6$. Bestimme $k$. \hfill (Abitur 2023 GK) \\ $k = $ \leerfeld
\teil Die Ebene $L$ enthält die x-Achse und den Punkt $G(3 | 2 | 5)$; ihre Gleichung hat die Form $r \cdot y + s \cdot z = 0$. Gib passende Werte für $r$ und $s$ an. \hfill (Abitur 2023 GK) \\ $r = $ \leerfeld , $s = $ \leerfeld
\teil Für jede Zahl $a$ ist $g_a: \vec{x} = (3 | 1 | -1) + t \cdot (1 | 2a | a)$. Jede Gerade $g_a$ liegt in der Ebene $E: y - 2z = c$. Bestimme $c$. \hfill (Abitur 2026 LK) \\ $c = $ \leerfeld
\teil Die Ebene $E: 2x + 3y - z = 8$ enthält einen Punkt, dessen drei Koordinaten übereinstimmen. Bestimme ihn. \hfill (Abitur 2019 LK) \\ \leerfeld
\end{teile}
\end{aufgabe}

%% TODO Titel: Ich-kann-Satz fehlt in der Bank (Platzhalter: Kettenname) – Nr. 13
%% TODO Anweisung: ein Satz über der ersten Teilaufgabe fehlt in der Bank – Nr. 13
\begin{aufgabe}{Parameter aus der Lagebedingung – Fehler finden, Begründen}
\begin{teile}
\teil Ich finde den Fehler: $E_k: kx + y + 2z = 8$ soll $P(2 | 4 | 1)$ enthalten. Jana setzt $Q(1 | 2 | 2)$ ein: $k + 2 + 4 = 8$, also $k = 2$.
\teil Begründe, warum das Einsetzen eines Punkts $P$ von $E_k: kx + 2y + z = 4$ eine Gleichung für $k$ liefert.
\end{teile}
\end{aufgabe}
